\documentclass[border=8pt]{standalone}
\usepackage{fontspec}
\setmainfont{Noto Sans}[Ligatures=TeXOff]
\XeTeXgenerateactualtext=1
\setsansfont{Noto Sans}
\newfontfamily\bengalifont{Noto Sans Bengali}[Script=Bengali, NFSSFamily=bnfam]
\usepackage[Latin,Bengali]{ucharclasses}
% group-free font switching (as in header-common.tex): robust inside braces, cells and quotes
\makeatletter
\def\BN@saved{\rmdefault}
\def\BN@bn{bnfam}
\newcommand\BNin{\edef\BN@tmp{\f@family}\ifx\BN@tmp\BN@bn\else\let\BN@saved\BN@tmp\fi\fontfamily{bnfam}\selectfont}
\newcommand\BNout{\fontfamily{\BN@saved}\selectfont}
\makeatother
\setTransitionsFor{Bengali}{\BNin}{\BNout}
% Typographic quotes and dashes are not in a ucharclasses class (class 0), so a Bangla word right
% after “ or ‘ would not trigger the Bangla font. Treat them as Latin so that every switch fires.
\XeTeXcharclass"2018=\XeTeXcharclass`A
\XeTeXcharclass"2019=\XeTeXcharclass`A
\XeTeXcharclass"201C=\XeTeXcharclass`A
\XeTeXcharclass"201D=\XeTeXcharclass`A
\XeTeXcharclass"2013=\XeTeXcharclass`A
\XeTeXcharclass"2014=\XeTeXcharclass`A
\XeTeXcharclass"2026=\XeTeXcharclass`A
\usepackage{newunicodechar}
\newunicodechar{।}{{\bengalifont\char"0964}}
\usepackage{tikz}
\usetikzlibrary{arrows.meta,positioning,shapes.geometric,shapes.misc,shapes.symbols,fit,backgrounds,calc,chains}
\definecolor{ink}{HTML}{111827}
\definecolor{band}{HTML}{F3F4F6}
\definecolor{band2}{HTML}{E5E7EB}
\definecolor{accent}{HTML}{1D4ED8}
\definecolor{accentlight}{HTML}{DBEAFE}
\definecolor{warn}{HTML}{B91C1C}
\definecolor{warnlight}{HTML}{FEE2E2}
\definecolor{ok}{HTML}{047857}
\definecolor{oklight}{HTML}{D1FAE5}
\tikzset{
  >={Stealth[length=2.6mm]},
  box/.style={draw=ink, line width=0.6pt, rounded corners=9pt, fill=white, align=center,
              inner sep=4pt, font=\small},
  pill/.style={box, rounded corners=13pt},
  key/.style={box, fill=accentlight, draw=accent, line width=0.9pt},
  soft/.style={box, fill=band},
  dec/.style={draw=ink, line width=0.6pt, diamond, aspect=2.2, align=center, inner sep=1.5pt,
              font=\footnotesize, fill=white},
  arr/.style={->, draw=ink, line width=0.6pt},
  darr/.style={->, draw=ink, line width=0.6pt, dashed},
  lbl/.style={font=\footnotesize, fill=white, inner sep=1.5pt},
  layer/.style={draw=band2, fill=band, rounded corners=6pt},
  title/.style={font=\bfseries\small, align=left, anchor=west},
}
\begin{document}
\sffamily\hyphenpenalty=10000\exhyphenpenalty=10000

\begin{tikzpicture}
\def\W{150mm}\def\H{104mm}
\fill[band] (0,0) rectangle (0.5*\W,0.5*\H);
\fill[accentlight!60] (0.5*\W,0) rectangle (\W,0.5*\H);
\fill[band2] (0,0.5*\H) rectangle (0.5*\W,\H);
\fill[accentlight] (0.5*\W,0.5*\H) rectangle (\W,\H);
\draw[ink, line width=0.6pt] (0,0) rectangle (\W,\H);
\draw[ink!50, dashed] (0.5*\W,0) -- (0.5*\W,\H);
\draw[ink!50, dashed] (0,0.5*\H) -- (\W,0.5*\H);
\draw[arr] (0,-6mm) -- node[below=1mm, font=\small]{BAPIL-এর কাজে আগ্রহ} (\W,-6mm);
\draw[arr] (-6mm,0) -- node[above=1mm, font=\small, rotate=90]{BAPIL-এর সাফল্যে প্রভাব} (-6mm,\H);
\node[font=\small\bfseries, anchor=north west, text=ink!70] at (2mm,\H-2mm) {সন্তুষ্ট রাখা};
\node[font=\small\bfseries, anchor=north east, text=accent] at (\W-2mm,\H-2mm) {ঘনিষ্ঠভাবে সম্পৃক্ত রাখা};
\node[font=\small\bfseries, anchor=south west, text=ink!70] at (2mm,2mm) {পর্যবেক্ষণ};
\node[font=\small\bfseries, anchor=south east, text=accent] at (\W-2mm,2mm) {অবহিত ও সম্পৃক্ত রাখা};
\tikzset{sh/.style={box, rounded corners=7pt, font=\footnotesize, inner sep=2.5pt}, shk/.style={sh, key}}
\node[shk] at (120mm,90mm) {আইসিটি বিভাগ};
\node[shk] at (104mm,78mm) {পরিকল্পনা কমিশন};
\node[shk] at (132mm,73mm) {অর্থ বিভাগ};
\node[sh] at (36mm,86mm) {মন্ত্রিপরিষদ বিভাগ};
\node[sh] at (44mm,70mm) {উপাত্ত ব্যবস্থাপনা কর্তৃপক্ষ};
\node[shk] at (98mm,62mm) {লেজিসলেটিভ বিভাগ};
\node[sh] at (128mm,58mm) {বিবিএস};
\node[sh] at (90mm,40mm) {এটুআই ও বিসিসি};
\node[sh] at (36mm,58mm) {উন্নয়ন সহযোগী};
\node[sh] at (130mm,42mm) {বিশ্ববিদ্যালয়};
\node[sh] at (110mm,26mm) {নাগরিক সমাজ};
\node[sh] at (60mm,36mm) {এআই হাব};
\node[sh] at (38mm,22mm) {শিল্প ও স্টার্টআপ};
\node[sh] at (124mm,12mm) {নাগরিক};
\node[sh] at (70mm,14mm) {সংবাদমাধ্যম};
\end{tikzpicture}
\end{document}
